\cleardoublepage
\section{Summary and Tips}

\subsection*{Interesting Facts}

\begin{itemize}
  \item The word \textbf{"algebra"} comes from the Arabic term \textbf{"al-jabr,"} which means "reunion of broken parts." It appeared in a famous math book written over 1,200 years ago by a scholar named Al-Khwarizmi. Many English math words like "algorithm" also come from his name!
  \item Even NASA engineers and video game developers simplify expressions when building models and writing code. It's a skill that goes far beyond school!
  \item Solving equations is a lot like solving a riddle or puzzle. You're looking for the one number that makes the equation true.
  \item People use multi-step equations in budgeting, engineering, and even cooking recipes when they need to scale quantities.
  \item Inequalities are used in many jobs—like making decisions in business, writing computer programs, or planning safe speeds in transportation.
\end{itemize}

\subsection*{Tips for Students}

\subsubsection*{Variables and Algebraic Expressions}
\begin{itemize}
  \item Think of a variable as a mystery box. What number could fit inside?
  \item Try different values to see how expressions change.
  \item When writing expressions, use colors, arrows, or circles to match terms and operations.
\end{itemize}

\subsubsection*{Simplifying Expressions}
\begin{itemize}
  \item Circle or highlight like terms before you combine them.
  \item Be careful with minus signs—they belong to the number or variable that follows!
  \item Go one step at a time to avoid careless mistakes.
\end{itemize}

\subsubsection*{Solving One-Step Equations}
\begin{itemize}
  \item Ask: "What's happening to the variable?" Then do the opposite.
  \item Keep both sides of the equation equal—whatever you do to one side, do to the other.
  \item Always plug your answer back in to make sure it works!
\end{itemize}

\subsubsection*{Solving Multi-Step Equations}
\begin{itemize}
  \item Don't rush—write each step clearly on its own line.
  \item Make sure each side of the equation is simplified before you solve.
  \item Always check your solution by plugging it back into the original equation.
\end{itemize}

\subsubsection*{Inequalities}
\begin{itemize}
  \item Always treat inequalities like equations—unless you multiply or divide by a negative!
  \item Graph solutions using number lines. Use:
    \begin{itemize}
      \item \textbf{Open circles} for $<$ or $>$
      \item \textbf{Closed circles} for $\leq$ or $\geq$
    \end{itemize}
  \item Check your solution by testing a number that fits the inequality.
\end{itemize}